\documentclass[10pt]{article}
\usepackage[margin=0.8in]{geometry}
\usepackage[T1]{fontenc}
\usepackage{amsmath,amssymb,booktabs,longtable,hyperref}
\hypersetup{colorlinks=true,urlcolor=blue,linkcolor=blue}
\title{Assessing Zhai et al. (2024) for Flow-Initialized RTI}
\author{Collision Avoidance Research}
\date{October 2, 2026}
\begin{document}
\maketitle

\section{Decision and scope}
The useful transfer from Zhai et al.~\cite{zhai} is motion-aware, coupled
path and speed guidance, rather than replacing the current initialization
with their complete artificial potential field (APF) controller.
The current nonlinear bicycle seeds already avoid sampled overlap in all
fourteen initial threat cases. Their more evident weaknesses are excessive
displacement, an approximate arrival clock, and an abrupt transition from
avoidance guidance to a free-pose terminal continuation.
These are concrete targets for improvement in an initializer used for one
local convex correction.

This assessment reads the supplied local paper, audits current code, measures
the unoptimized seeds, and performs one controlled ablation. It does not
implement the paper's complete algorithm or replace the production controller.
The proposed predictive-field design below is a research decision, not an
experimentally established improvement or a safety certificate.
Only the diagnostic driver and research artifacts are new.

\section{What the paper actually contributes}
The primary source is the supplied 17-page IEEE Access paper, volume 12,
pp.\ 19665--19681, DOI 10.1109/ACCESS.2024.3355952. The local PDF was read as text;
the equations on PDF page 7 were also inspected visually. Crossref's publisher
metadata confirms title, authors, year, volume and pages. IEEE web access was
challenge-blocked, so substantive claims come from the local paper.

\begin{itemize}
\item Section III-B, Eqs.\ (20)--(22), makes repulsion angle dependent through
  a water-droplet-shaped field. Eqs.\ (25)--(30) add terms involving relative
  velocity and acceleration for moving obstacles.
\item Section IV-B/C explicitly reasons about collision-area entry/exit times
  and chooses speed changes. Table 1 defines eight prioritized driving
  behaviors; Figs.\ 17--18 contain overtaking/crossing decision logic.
  These are part of the reported method, not properties of a potential field
  alone.
\item Section V-A adds stepwise speed, heading and acceleration restrictions
  and states that a Bezier curve smooths the planned trajectory. A tracked
  vehicle model and a separate tracking controller execute it.
\item In the lateral-meeting experiment, Figs.\ 25--28, APF without the speed
  planner collides; with speed planning, the vehicle slows and passes safely.
  This directly supports the relevance of arrival timing, but does not isolate
  every contribution of the field, behavior rules, smoother and tracker.
\end{itemize}

The principal integrated simulations use a target cruise speed of 18 km/h
(5 m/s); the earlier qualitative path comparison in Fig.\ 13 uses a 10 m/s
ego. Neither is our 15 m/s Fiala-bicycle RTI initialization experiment.
The real-vehicle demonstration in Section V-D concerns static obstacles.
No reproducible per-frame computation-time distribution was found in the
paper. Its HIL and tracking results do not establish a 100 ms bound for our
controller or prove that one convex correction will succeed.

\paragraph{What should not transfer directly.}
The eight-behavior arbitration, virtual point destinations, tracked-vehicle
drive/braking equations, and Eq.\ (41)'s stationary endpoint conditions are not
needed for our single-controller, given-path cruise task. A smooth geometric
Bezier curve is not automatically a trajectory of our nonlinear bicycle.
Any geometric or speed guidance must still be converted to inputs and rolled
out with the same bicycle model used for linearization.

\section{Current implementation and measured weaknesses}
The source baseline is commit
\path{eab55398f9a880666fdb5b078b41a77b291b136e}.

The audit examines the following routines:
\begin{itemize}
\item \path{controller/predictiveSafetyGeometry.m}: the Gaussian guide
and transported flow velocity.
\item \path{controller/solvePredictiveControl.m}: the nonlinear flow seed,
input shaping and clipping.
\end{itemize}
Normal warm operation shifts previous inputs; flow initialization is used
when a usable sequence is absent or a fresh seed is requested. It is not
regenerated on every successful warm frame.

\paragraph{Motion is already included.}
The target follows the exact constant-tangential-acceleration,
constant-sideslip model at matching absolute times. A Gaussian envelope is fit
using the approximate ego station clock
\[
 \bar s(t)=s_0+v_{\rm ref}t,\qquad
 d_g(t)=A_g\exp[-(t-t_c)^2/(2\sigma^2)].
\]
That guidance is then modified by a transported cylinder flow, mapped to
heading/speed requests, shaped using tire capacity, and rolled out with the
nonlinear bicycle. Thus it would be incorrect to diagnose the current
initializer as ignoring target motion.

\paragraph{Four specific issues.}
\begin{enumerate}
\item The Gaussian has no initial-position or initial-tangent matching.
  Its requested lateral offset at $t=0$ is 0.184--4.765 m in these cases,
  although the actual ego starts on its given path. The nonlinear rollout
  starts at the correct measured state; the discontinuity is in the guidance
  request, not a fabricated state jump.
\item The shape is a circular exclusion envelope of radius 5.262 m for these
  vehicles, independent of their relative orientation. For parallel,
  side-by-side 4.8 m by 1.9 m rectangles, a 0.1 m clearance requires only
  2.0 m of lateral center separation. This comparison illustrates the
  initializer's geometric conservatism; it is not a clearance calculation
  for the actual turning motion. The optimized collision rows already use
  rectangle geometry and should retain it.
\item The clock used to fit the Gaussian is not updated when the generated
  ego turns or slows. Actual seed station differs from that clock by up to
  6.789 m. Target prediction remains time aligned; the inaccurate part is
  the anticipated ego arrival. Also, clipping the norm of the flow velocity
  to a desired speed is not an explicit longitudinal collision-timing plan.
\item Flow guidance ends at the last overlap of the nominal station window,
  not after its Gaussian has decayed. At the twelve actual handoffs,
  its hypothetical remaining offset is 5.121--7.427 m. The seed then follows
  a straight terminal family whose pose is freely refitted. This settles
  intrinsic dynamics but can preserve an outward heading. The original
  input changes by about 0.175--0.204 rad at that handoff. A free terminal
  pose is not itself a defect; choosing its approach this abruptly biases
  the initialization away from nominal recovery.
\end{enumerate}

\subsection{Focused diagnostic experiment}
\path{scripts/auditFlowInitialization.m} extracts the current seed,
shaping and clipping helpers verbatim into the output directory. It evaluates
seven existing threat scenarios at 8 and 15 m/s, with 50 ms holds, prefixes
of 8 and 16 holds, and the unchanged three-second completion. The resulting
3.4/3.8 s lengths are the current initializer lengths, not success deadlines
for a closed-loop experiment.
States are exact, road constraints absent, and no random seed is used.

The audit independently integrates the whole prescribed input sequence with
ODE45, relative tolerance $10^{-10}$, absolute tolerance $10^{-11}$, and
21 output points per hold. It compares rectangles at matching times.
The table gives minimum sampled physical clearance and maximum lateral
displacement along the original RK4 anchor.
No sampled overlaps occur in these initial seeds; this is not an all-time
certificate or a test of flow restarts at arbitrary later states.

The ablation changes only post-flow guidance from the free-pose terminal
feedback to the existing nominal path guidance. Every preceding input is
asserted identical. It is a diagnostic copy, not a controller option.

\begin{center}
\small
\begin{tabular}{rlrrrr}
\toprule
$v_{\rm ref}$ & Scenario & Gap (m) & Original $|e_y|_{\max}$ (m)
 & Ablation $|e_y|_{\max}$ (m) & Original core ratio\\
\midrule
8 & Head-on & 2.184 & 14.928 & 9.310 & 0.0100\\
8 & Accelerating head-on & 1.757 & 15.711 & 9.327 & 0.0029\\
8 & Braking lead & 1.293 & 6.807 & 6.807 & 4408.6\\
8 & Crossing & 0.860 & 12.110 & 9.684 & 0.7900\\
8 & Turning crossing & 1.133 & 12.120 & 10.117 & 1.5432\\
8 & Curved head-on & 2.678 & 14.055 & 9.292 & 0.0047\\
8 & Curved crossing & 0.698 & 11.544 & 9.566 & 0.8708\\
15 & Head-on & 0.288 & 15.950 & 6.782 & 0.000054\\
15 & Accelerating head-on & 0.172 & 15.345 & 6.260 & 0.000037\\
15 & Braking lead & 1.284 & 8.504 & 8.504 & 654.0\\
15 & Crossing & 3.123 & 19.269 & 12.969 & 0.0099\\
15 & Turning crossing & 3.154 & 19.983 & 13.238 & 0.0172\\
15 & Curved head-on & 3.327 & 16.907 & 10.792 & 0.0033\\
15 & Curved crossing & 2.310 & 15.551 & 10.517 & 0.0640\\
\bottomrule
\end{tabular}
\end{center}

All ablation minimum gaps equal the corresponding original values to the
reported precision; the minima precede the changed tail. Maximum displacement
falls in twelve cases; the two braking cases never enter the changed branch.
This isolates a contribution of the handoff to the large displacement.
It does not establish a superior initializer: the terminal-core result worsens.

The core ratio is
\[
 \eta_T=\frac{\|R_q([v_x,v_y,r,u_{H-1}^{\top}]^\top
           -[v_{x,f},v_{y,f},r_f,u_f^\top]^\top)\|_2}{\rho_f}.
\]
It minimizes out the free endpoint pose; $\eta_T\leq1$ means intrinsic
membership only, without asserting terminal target separation. Eleven original
anchors satisfy this condition, whereas none of the fourteen ablation anchors
does. For example, 15 m/s head-on improves displacement from 15.950 to
6.782 m but changes $\eta_T$ from $5.36\,10^{-5}$ to 471.6.
Forcing an immediate return turn trades settled endpoint dynamics for a
prettier path. It is not a valid drop-in solution to the RTI problem.
The two original braking seeds and the 8 m/s turning-crossing seed already
need endpoint repair; this does not imply the closed-loop controller fails.

Seven alternating warm calls per helper give a median across case medians
of 3.351 ms for the current seed and 3.250 ms for the ablation. These exclude
terminal preparation, optimization and the audit. The largest RK4-versus-ODE45
body-pose discrepancy across an entire seed is 7.07 mm for the current seeds
and 9.91 mm for the ablation. These are numerical comparisons, not model-error
bounds. No new closed-loop recovery or controller latency claim is made.

\section{Printed equations require interpretation before implementation}
The following are checks of the printed paper, not claims about the authors'
unavailable executable implementation.

\begin{enumerate}
\item Eq.\ (23) is not a consistent translated rotation as printed:
  substituting $p=p_{\rm obs}=(3,4)$ and $\theta=\pi/6$ gives $(-4,3)$,
  rather than zero. A relative frame must use one consistent rotation of
  $p-p_{\rm obs}$.
\item Even holding $k_d$ constant, differentiating Eq.\ (20) gives radial
  force $k_{\rm rep}(1/(k_dd)-1/d_0)/(k_dd^2)$.
  Eq.\ (22) prints $k_d^2$ in the prefactor denominator. With $k_d=1.2$,
  $d=4$, $d_0=10$, and $k_{\rm rep}=1$, the negative finite-difference
  derivative is 0.00564236; the printed force is 0.00470197.
  Since $k_d$ actually depends on angle, a full spatial gradient also has
  angular terms.
\item The printed $d\leq d_0$ cutoff is not consistent with zeroing a potential
  based on $k_dd$ at the boundary. For example, $k_d=1.6$, $d=8$, $d_0=10$
  makes the radial expression in Eq.\ (22) negative. The cutoff,
  effective distance and gradient should be derived together if this
  field is adapted.
\end{enumerate}

These issues, the behavior-specific decision logic, and the absence of a
complete reusable parameter/smoothing specification make literal reproduction
in our controller inappropriate. They do not invalidate the experimental
observation that coordinated speed planning helped the paper's crossings.

\section{A compatible design direction}
Retain one initializer and the same PCBF--CLF controller. Use a predictive
obstacle field to guide both course and speed, with a circulation component
to resolve symmetric head-on encounters. Keep the exact target trajectory
and nonlinear bicycle rollout. The flow remains a search mechanism, not an
executable fallback or a second CLF.

One suitable design template is a fixed-size, time-synchronized risk function
\[
 {\cal R}(p,v,\chi,t)
   =\sum_{j=1}^{Q}w_j\,
     \Phi\!\left(D\big(\widehat{\cal E}(p,v,\chi;\tau_j),
                                {\cal T}(t+\tau_j)\big)\right).
\]
Here $Q$ is bounded, $\widehat{\cal E}$ is a cheap local ego-motion predictor,
and $D$ is an orientation-aware signed separation surrogate used only for
guidance. It should retain a useful direction in overlap, unlike a
nonnegative ordinary distance clipped to zero. A smooth repulsion
$\Phi$ increases as separation decreases. The current circle may be
replaced by a directional rectangle/ellipse approximation without treating
that approximation as the controller's safety geometry.
The course and speed enter the same future collision prediction, instead of
fitting a lateral curve under a permanently frozen cruise clock.

For a local predictor
$\hat p_E(\tau)=p+\tau v(\cos\chi,\sin\chi)^\top$, its analytic sensitivities are
\[
 \partial_v\hat p_E=\tau(\cos\chi,\sin\chi)^\top,\qquad
 \partial_\chi\hat p_E=\tau v(-\sin\chi,\cos\chi)^\top.
\]
They yield inexpensive joint guidance, for example
\[
 \begin{bmatrix}a_g\\\omega_g\end{bmatrix}
 =
 \begin{bmatrix}a_{\rm path}\\\omega_{\rm path}\end{bmatrix}
 -
 \begin{bmatrix}g_v&0\\0&g_\chi\end{bmatrix}
 \begin{bmatrix}\partial_v{\cal R}\\\partial_\chi{\cal R}\end{bmatrix}
 +\begin{bmatrix}0\\\omega_{\rm circulation}\end{bmatrix}.
\]
The gains carry the necessary units or act on normalized variables.
This is an explicit guidance update, not another numerical optimization.
The baseline path field attracts the vehicle toward the original path and
speed; risk modulates longitudinal progress and turning together.
Circulation preserves a consistent passing side when a symmetric gradient
alone cannot choose one.

At each seed hold, map this guidance to tire-informed inputs, advance
$\bar x_{i+1}=f_h(\bar x_i,\bar u_i)$ from the measured state, and recompute
risk at that new state and absolute time. The virtual predictor is only
guidance; it does not replace the bicycle dynamics in the anchor.
Smooth start, obstacle-field decay and endpoint settling must be designed
together. In particular, the terminal pose remains free: there should be no
requirement to reach a fixed road position at a fixed time, nor to complete
nominal recovery within the short seed horizon.

This template transfers the paper's most relevant insight while avoiding its
behavior library and printed-formula ambiguities. Its open questions are
the surrogate shape, circulation strength, normalized gains and endpoint
settling. They need a focused implementation comparison before adoption.
There is no claim that APF eliminates local minima or that a better seed
alone guarantees one-solve nonlinear feasibility.

\paragraph{Acceptance criteria for that future implementation.}
Compare raw matching-time clearance, required input correction, endpoint
repair, and first-convex-correction model agreement, not just path length.
Then use all fourteen complete closed loops through nominal recovery and
measure total call latency, including initialization and exceptional rebuilds.
A desirable initializer should improve these quantities together.
The present tail ablation fails the endpoint criterion, despite reducing
displacement, and is therefore not promoted into production.

\section{Reproduction and evidence}
Run from the repository root:
\begin{quote}\small
\texttt{matlab -singleCompThread -batch}\\
\texttt{"addpath('scripts'); auditFlowInitialization('/tmp/flow-audit');"}
\end{quote}
The driver emits the metrics, full traces and identified extracted diagnostic
helpers. Compact committed evidence is under
\path{report/FLOW_INITIALIZATION_ZHAI_20261002/}.
It includes source/artifact provenance, the final 28-row audit and an
independent numerical check of the printed equations.
The report uses tool-assisted PDF extraction, code comparison and numerical
analysis; source claims, measured findings and proposed adaptations are
separated explicitly.

\begin{thebibliography}{1}
\bibitem{zhai}
Zhai, L., Liu, C., Zhang, X., \& Wang, C. (2024).
Local trajectory planning for obstacle avoidance of unmanned tracked vehicles
based on artificial potential field method.
\emph{IEEE Access, 12}, 19665--19681.
\href{https://doi.org/10.1109/ACCESS.2024.3355952}{doi:10.1109/ACCESS.2024.3355952}.
\end{thebibliography}
\end{document}
